\section{Conclusion}
\label{sec:conclusion}

We investigate whether the attention mechanism is a universal computational principle that extends beyond transformers and NLP. Through cross-domain experiments, attention pattern analysis, and formal mathematical mapping, we find strong evidence supporting this view. Attention-based models outperform non-attention baselines on text from scratch ($p{=}0.046$) and dominate across both text and vision at scale with pretraining (\vit outperforms \resnet by 22 percentage points). Attention patterns adapt to domain structure---diffuse for text, focused for vision---without any architectural modification. Most importantly, we formally show that transformer self-attention, PageRank, and collaborative filtering are instances of the same mathematical operation: normalized relevance-weighted aggregation.

The key takeaway is that attention---understood as selective information routing through a capacity bottleneck---is not a transformer-specific technique but a unifying principle across AI architectures and internet systems. This perspective suggests that attention-based architectures should be the default starting point for new domains and tasks, and that the mathematical connection between computational and economic attention deserves deeper theoretical investigation.

Future work should extend the empirical evaluation to audio, graph, and multimodal domains; develop a quantitative (not just structural) theory linking attention mechanisms to user engagement in recommendation systems; and investigate whether sub-quadratic alternatives to attention~\citep{gu2023mamba} preserve the universality properties we identify or sacrifice them for efficiency.
